% !PRINTED 439
\weberpage{439}
%<<CONT
are connected (\S.\,37, 9.). Thus by $A$ not only is a point derived from every point, but also a straight line from every straight line.

The relation of the points $y$ to the points $x$ expressed by (1) can be designated as a \gesperrt{mapping of the plane into itself}, in which to every point there corresponds a definite other point, and to every straight line a straight line. Such a mapping is also called in geometry a collineation. The point $y$ is called the \gesperrt{image} of the point $x$, and the straight line $v$ the \gesperrt{image} of the straight line $u$. If one has the images of two points, then the connecting line of these images is the image of the connecting line of the two original points.

If $S$ denotes a second ternary linear substitution and
\[ A' = S^{-1}AS \]
the substitution transformed from $A$ by $S$, then
\[ (y') = A'(x') \]
is equivalent to
\[ (y) = A(x), \]
if
\[ (y) = S(y'), \quad (x) = S(x') \]
is set. Now instead of defining a mapping by $S$, one can also understand by it a \gesperrt{coordinate transformation}, by understanding by $x_1', x_2', x_3'$ the coordinates of the point $(x)$ in a new coordinate system determined by $S$. Then the connection between the two points $x$, $y$ is expressed by the transformed substitution $A'$ just as by $A$, and the transformation of the substitutions is thus equivalent to the transformation of the coordinate system. The composition of the transformed substitutions takes place, as we have already seen earlier, just like that of the original ones, and thus by transformation there arises from every group an isomorphic group.

Of great importance is now the investigation of the points that remain unchanged under the mapping by a substitution $A$.

Let us take $A$ in the form:
\[ A = \begin{pmatrix} \alpha, & \beta, & \gamma \\ \alpha', & \beta', & \gamma' \\ \alpha'', & \beta'', & \gamma'' \end{pmatrix}, \]
%CONT>>
